\documentclass[10pt,a4paper]{amsart}
\usepackage{amsmath,amssymb,amsthm,graphicx}
\usepackage{verbatim}
\newenvironment{reference}{\comment}{\endcomment}
%\newenvironment{reference}{}{}
\newenvironment{slogan}{\comment}{\endcomment}
\newenvironment{history}{\comment}{\endcomment}

% For commutative diagrams we use Xy-pic
\usepackage[all]{xy}

% We use 2cell for 2-commutative diagrams.
\xyoption{2cell}
% Theorem environments.
%
\theoremstyle{plain}
\newtheorem{theorem}[subsection]{Theorem}
\newtheorem{proposition}[subsection]{Proposition}
\newtheorem{lemma}[subsection]{Lemma}

\theoremstyle{definition}
\newtheorem{definition}[subsection]{Definition}
\newtheorem{example}[subsection]{Example}
\newtheorem{exercise}[subsection]{Exercise}
\newtheorem{situation}[subsection]{Situation}

\theoremstyle{remark}
\newtheorem{remark}[subsection]{Remark}
\newtheorem{remarks}[subsection]{Remarks}

\numberwithin{equation}{subsection}

% Macros
%
\def\lim{\mathop{\mathrm{lim}}\nolimits}
\def\colim{\mathop{\mathrm{colim}}\nolimits}
\def\Spec{\mathop{\mathrm{Spec}}}
\def\Hom{\mathop{\mathrm{Hom}}\nolimits}
\def\Ext{\mathop{\mathrm{Ext}}\nolimits}
\def\SheafHom{\mathop{\mathcal{H}\!\mathit{om}}\nolimits}
\def\SheafExt{\mathop{\mathcal{E}\!\mathit{xt}}\nolimits}
\def\Sch{\mathit{Sch}}
\def\Mor{\mathop{\mathrm{Mor}}\nolimits}
\def\Ob{\mathop{\mathrm{Ob}}\nolimits}
\def\Sh{\mathop{\mathit{Sh}}\nolimits}
\def\NL{\mathop{N\!L}\nolimits}
\def\CH{\mathop{\mathrm{CH}}\nolimits}
\def\proetale{{pro\text{-}\acute{e}tale}}
\def\etale{{\acute{e}tale}}
\def\QCoh{\mathit{QCoh}}
\def\Ker{\mathop{\mathrm{Ker}}}
\def\Im{\mathop{\mathrm{Im}}}
\def\Coker{\mathop{\mathrm{Coker}}}
\def\Coim{\mathop{\mathrm{Coim}}}

% Boxtimes
%
\DeclareMathSymbol{\boxtimes}{\mathbin}{AMSa}{"02}

%
% Macros for moduli stacks/spaces
%
\def\QCohstack{\mathcal{QC}\!\mathit{oh}}
\def\Cohstack{\mathcal{C}\!\mathit{oh}}
\def\Spacesstack{\mathcal{S}\!\mathit{paces}}
\def\Quotfunctor{\mathrm{Quot}}
\def\Hilbfunctor{\mathrm{Hilb}}
\def\Curvesstack{\mathcal{C}\!\mathit{urves}}
\def\Polarizedstack{\mathcal{P}\!\mathit{olarized}}
\def\Complexesstack{\mathcal{C}\!\mathit{omplexes}}
% \Pic is the operator that assigns to X its picard group, usage \Pic(X)
% \Picardstack_{X/B} denotes the Picard stack of X over B
% \Picardfunctor_{X/B} denotes the Picard functor of X over B
\def\Pic{\mathop{\mathrm{Pic}}\nolimits}
\def\Picardstack{\mathcal{P}\!\mathit{ic}}
\def\Picardfunctor{\mathrm{Pic}}
\def\Deformationcategory{\mathcal{D}\!\mathit{ef}}
\usepackage{cleveref}
\crefname{originalreference}{original reference}{original references}
\Crefname{originalreference}{Original reference}{Original references}
\makeatletter
\newlabel{lemma-algebraic-permanence}{{09GJ}{source}{}{sourcepointer-1}{}}
\newlabel{lemma-algebraic-permanence@cref}{{[originalreference][0][]09GJ}{[1][0][]source}{}{}{}}
\makeatother
\crefname{originalreference}{original reference}{original references}
\Crefname{originalreference}{Original reference}{Original references}
\makeatletter
\newlabel{lemma-algebraically-closed}{{09GR}{source}{}{sourcepointer-2}{}}
\newlabel{lemma-algebraically-closed@cref}{{[originalreference][0][]09GR}{[1][0][]source}{}{}{}}
\makeatother
\crefname{originalreference}{original reference}{original references}
\Crefname{originalreference}{Original reference}{Original references}
\makeatletter
\newlabel{lemma-finite-is-algebraic}{{09GG}{source}{}{sourcepointer-3}{}}
\newlabel{lemma-finite-is-algebraic@cref}{{[originalreference][0][]09GG}{[1][0][]source}{}{}{}}
\makeatother
\crefname{originalreference}{original reference}{original references}
\Crefname{originalreference}{Original reference}{Original references}
\makeatletter
\newlabel{lemma-size-algebraic-extension}{{09GK}{source}{}{sourcepointer-4}{}}
\newlabel{lemma-size-algebraic-extension@cref}{{[originalreference][0][]09GK}{[1][0][]source}{}{}{}}
\makeatother
\title{Every field has an algebraic closure}
\author{The Stacks Project Authors}
\date{stacks-project:fields.tex:local-snapshot-20261012; GNU FDL 1.2 or later}
\begin{document}
\maketitle
\noindent\textit{Editorial note.} All delivered mathematical text is the author's own, copied from the pinned Stacks Project chapter file fields.tex; no statement or proof step was rewritten, reordered or removed.\par\medskip
\noindent\textit{Editorial note.} The counted proof is the author's complete argument: The author interjects before the proof (source lines 1188-1191) that this existence result is mostly a red herring for the rest of the chapter; the interjection is outside the counted span and is not delivered.\par\medskip
\noindent\textit{Editorial note.} The label theorem-existence-algebraic-closure names the counted statement, and the proof environment immediately following it in the source is delivered verbatim as its proof.\par\medskip
\noindent\textit{Editorial note.} The counted statement is carried by the author's own theorem environment in the chapter source.\par\medskip
\noindent\textit{Allowed standard prerequisites (editorial).}\begin{itemize}
\item Context span 1169-1177: definition of algebraic closure (algebraic + algebraically closed)
\item Same chapter: lemma-size-algebraic-extension (cardinality of an algebraic extension of F is at most max(aleph\_0, |F|)), lemma-algebraically-closed (equivalent characterizations, so not algebraically closed means an irreducible polynomial of degree > 1 exists), lemma-finite-is-algebraic and lemma-algebraic-permanence (algebraicity is stable in towers)
\end{itemize}\medskip
\ifdefined\hypertarget\hypertarget{sourcepointer-1}{}\fi
\noindent\textit{Source reference.} Original reference 09GJ: sources/stacks-fields/pinned-original-source.tex, source lines 957--957. Prerequisite pointer invoked at 1231; it names a standard fact of the same chapter that the counted argument uses, and every step of the counted proof built on it is retained verbatim in the delivered body.\par\medskip
\ifdefined\hypertarget\hypertarget{sourcepointer-2}{}\fi
\noindent\textit{Source reference.} Original reference 09GR: sources/stacks-fields/pinned-original-source.tex, source lines 1128--1128. Prerequisite pointer invoked at 1228; it names a standard fact of the same chapter that the counted argument uses, and every step of the counted proof built on it is retained verbatim in the delivered body.\par\medskip
\ifdefined\hypertarget\hypertarget{sourcepointer-3}{}\fi
\noindent\textit{Source reference.} Original reference 09GG: sources/stacks-fields/pinned-original-source.tex, source lines 878--878. Prerequisite pointer invoked at 1229; it names a standard fact of the same chapter that the counted argument uses, and every step of the counted proof built on it is retained verbatim in the delivered body.\par\medskip
\ifdefined\hypertarget\hypertarget{sourcepointer-4}{}\fi
\noindent\textit{Source reference.} Original reference 09GK: sources/stacks-fields/pinned-original-source.tex, source lines 986--986. Prerequisite pointer invoked at 1194; it names a standard fact of the same chapter that the counted argument uses, and every step of the counted proof built on it is retained verbatim in the delivered body.\par\medskip
\begin{definition}
\label{definition-algebraic-closure}
Let $F$ be a field. An {\it algebraic closure} of $F$ is a field
$\overline{F}$ containing $F$ such that:
\begin{enumerate}
\item $\overline{F}$ is algebraic over $F$.
\item $\overline{F}$ is algebraically closed.
\end{enumerate}
\end{definition}

\begin{theorem}
\label{theorem-existence-algebraic-closure}
Every field has an algebraic closure.
\end{theorem}

\begin{proof}
Let $F$ be a field. By Lemma \ref{lemma-size-algebraic-extension} the
cardinality of an algebraic extension of $F$ is bounded by
$\max(\aleph_0, |F|)$. Choose a set $S$ containing $F$ with
$|S| > \max(\aleph_0, |F|)$. Let's consider triples
$(E, \sigma_E, \mu_E)$ where
\begin{enumerate}
\item $E$ is a set with $F \subset E \subset S$, and
\item $\sigma_E : E \times E \to E$ and $\mu_E : E \times E \to E$
are maps of sets such that $(E, \sigma_E, \mu_E)$ defines the structure
of a field extension of $F$ (in particular $\sigma_E(a, b) = a +_F b$
for $a, b \in F$ and similarly for $\mu_E$), and
\item $E/F$ is an algebraic field extension.
\end{enumerate}
The collection of all triples $(E, \sigma_E, \mu_E)$ forms a set $I$.
For $i \in I$ we will denote $E_i = (E_i, \sigma_i, \mu_i)$ the
corresponding field extension to $F$. We define a partial ordering on
$I$ by declaring $i \leq i'$ if and only if $E_i \subset E_{i'}$
(this makes sense as $E_i$ and $E_{i'}$ are subsets of the same set $S$)
and we have $\sigma_i = \sigma_{i'}|_{E_i \times E_i}$ and
$\mu_i = \mu_{i'}|_{E_i \times E_i}$, in other words, $E_{i'}$ is a field
extension of $E_i$.

\medskip\noindent
Let $T \subset I$ be a totally ordered subset. Then it is clear that
$E_T = \bigcup_{i \in T} E_i$ with induced maps $\sigma_T = \bigcup \sigma_i$
and $\mu_T = \bigcup \mu_i$ is another element of $I$. In other words
every totally order subset of $I$ has a upper bound in $I$. By Zorn's lemma
there exists a maximal element $(E, \sigma_E, \mu_E)$ in $I$. We claim that
$E$ is an algebraic closure. Since by definition of $I$ the extension
$E/F$ is algebraic, it suffices to show that $E$ is algebraically closed.

\medskip\noindent
To see this we argue by contradiction. Namely, suppose that $E$ is not
algebraically closed. Then there exists an irreducible polynomial
$P$ over $E$ of degree $> 1$, see Lemma \ref{lemma-algebraically-closed}.
By Lemma \ref{lemma-finite-is-algebraic} we obtain a nontrivial finite
extension $E' = E[x]/(P)$. Observe that $E'/F$ is algebraic by
Lemma \ref{lemma-algebraic-permanence}.
Thus the cardinality of $E'$ is $\leq \max(\aleph_0, |F|)$.
By elementary set theory we can extend the given injection
$E \subset S$ to an injection $E' \to S$. In other words, we may
think of $E'$ as an element of our set $I$ contradicting the
maximality of $E$. This contradiction completes the proof.
\end{proof}

\end{document}
